\subsection{Operations on complexes}

On the set A × A, the two operations

\[
\begin{array}{l} (a, b) + (a ^ {\prime}, b ^ {\prime}) = (a + a ^ {\prime}, b + b ^ {\prime}) \\ (a, b) (a ^ {\prime}, b ^ {\prime}) = (a a ^ {\prime}, a b ^ {\prime} + b a ^ {\prime}) \end{array}
\]

define a ring structure, denoted A(ε), with identity element 1 = (1, 0); the injection a ↦ (a, 0) = a1 allows A to be identified with a subring of A(ε); the module A(ε) is free with basis \{ 1, ε \} where ε = (0, 1); we have ε² = 0 and ε is central in A(ε).

When A is commutative, \(\mathrm{A}(\varepsilon)\) is a dual number algebra over A (III, p.~15). Endow \(\mathrm{A}(\varepsilon)\) with the ring grading (II, p.~164) for which \(\mathrm{A}(\varepsilon)_{0} = \mathrm{A}.1\) , \(\mathrm{A}(\varepsilon)_{-1} = \mathrm{A}. \varepsilon\) and \(\mathrm{A}(\varepsilon)_{n} = 0\) for \(n \neq 0, -1\) . It is clear that giving a structure of A-complex on a set C amounts to giving on C a structure of \(\mathrm{A}(\varepsilon)\) graded module, the differential d corresponding to the homothety \(\varepsilon_{C}\) ; likewise, morphisms of complexes correspond to graded homomorphisms of degree 0 of \(\mathrm{A}(\varepsilon)\) graded modules. The species of structure of \(\mathrm{A}(\varepsilon)\) graded modules and of A-complexes are therefore equivalent (E, IV, p.~9-10). We shall use this fact to transport to the theory of complexes the usual notions of the theory of graded modules.

To the notion of graded sub-A(ε)-module corresponds that of subcomplex: a subcomplex of the complex (C, d) is therefore a graded submodule C' of C such that, for every \(n \in Z\) , \(d_{n}(C_{n}') \subset C_{n-1}'\) ; if we denote by \(d'\) the graded A-homomorphism from C' to C' induced by d, then (C', d') is a complex structure, called the structure induced by (C, d). Unless expressly stated otherwise, every subcomplex will be equipped with the induced structure.

We leave it to the reader to similarly spell out the notions of quotient complex, exact sequence of complexes, kernel, cokernel, and image of a morphism of complexes, following the dictionary below:

graded quotient A(ε)-module

= quotient complex,

kernel, cokernel, image of a graded A(ε)-homomorphism of degree 0 \} kernel, cokernel, image of a morphism of complexes,

exact sequence of graded A(ε)-modules and graded homomorphisms of degree 0 \} = exact sequence of complexes.

For example, the canonical exact sequences of No. 1 give exact sequences of complexes, called canonical.

\begin{enumerate}
\def\labelenumi{(\Roman{enumi})}
\tightlist
\item
\end{enumerate}

\[
0 \longrightarrow \mathrm{Z} (\mathrm{C}) \longrightarrow \mathrm{C} \stackrel {\delta} {\longrightarrow} \mathrm{B} (\mathrm{C}) (- 1) \longrightarrow 0,\tag{II}
\]

\[
0 \longrightarrow \mathrm{B} (\mathrm{C}) \longrightarrow \mathrm{Z} (\mathrm{C}) \longrightarrow \mathrm{H} (\mathrm{C}) \longrightarrow 0,\tag{III}
\]

\[
0 \longrightarrow \mathrm{B} (\mathrm{C}) \longrightarrow \mathrm{C} \longrightarrow \mathrm{C} / \mathrm{B} (\mathrm{C}) \longrightarrow 0,\tag{IV}
\]

\[
0 \longrightarrow \mathrm{H} (\dot {\mathrm{C}}) \longrightarrow \mathrm{C} / \mathrm{B} (\mathrm{C}) \stackrel {\delta} {\longrightarrow} \mathrm{B} (\mathrm{C}) (- 1) \stackrel {\cdot} {\longrightarrow} 0,\tag{V}
\]

\[
0 \longrightarrow \mathrm{H} (\mathrm{C}) \longrightarrow \mathrm{C} / \mathrm{B} (\mathrm{C}) \longrightarrow \mathrm{Z} (\mathrm{C}) (- 1) \longrightarrow \mathrm{H} (\mathrm{C}) (- 1) \longrightarrow 0.
\]

Similarly, one defines the notions of direct sum of complexes, inductive system of complexes, and inductive limit of an inductive system of complexes.

Let \((\mathbf{C}_{i}, d_{i})\) be a family of complexes. One calls product of this family, and one denotes by \(\prod_{i \in I} (\mathbf{C}_{i}, d_{i})\), the complex \((\mathbf{C}, d)\) obtained as follows:

\begin{enumerate}
\def\labelenumi{\alph{enumi})}
\item
  for each \(n \in Z\) , \(C_{n}\) is the product A-module \(\prod_{i \in I} (\mathbf{C}_{i})_{n}\) of the homogeneous components \((\mathbf{C}_{i})_{n}\) of the given complexes,
\item
  for each \(n \in Z\) , \(d_{n} : C_{n} \to C_{n-1}\) is the A-homomorphism with components \((d_{i})_{n}\) .
\end{enumerate}

When I is finite, \(\prod_{i\in I}(\mathbf{C}_{i},d_{i})\) is equal to \(\bigoplus_{i\in I}(\mathbf{C}_{i},d_{i})\) One should note that, in general, the underlying A-module of the product complex \(\prod_{i\in I}(\mathbf{C}_{i},d_{i})\) is not the product A-module \(\prod_{i\in I}\mathbf{C}_{i}\) .

Consider a family (resp. a filtered inductive system) \((\mathbf{C}_i)_{i\in I}\) Let C be the direct sum (resp. the inductive limit) of \(C_i\) and let \(\alpha_{i}:C_{i}\to C\) the canonical homomorphisms. Then the \(\mathrm{H}(\alpha_{i}): \mathrm{H}(\mathrm{C}_{i})\to \mathrm{H}(\mathrm{C})\) define a graded homomorphism of degree 0, called canonical, from \(\bigoplus_{i\in I}\mathrm{H}(\mathrm{C}_{i})\) (resp. \(\varinjlim_{i\in I}\mathrm{H}(\mathrm{C}_{i})\) ) in \(\mathrm{H}(\mathrm{C})\) . Likewise, the canonical projections \(\prod_{i\in I}\mathrm{C}_{i}\to \mathrm{C}_{i}\) define a graded homomorphism of degree 0, called canonical, from \(\mathrm{H}(\prod \mathrm{C}_i)\) to \(\prod (\mathrm{H}(\mathrm{C}_i))\) .

PROPOSITION 1. --- For every family of complexes \((\mathbf{C}_{i})_{i\in I}\) , the canonical homomorphisms

\[
\bigoplus_ {i \in \mathbf {I}} \mathrm{H} (\mathrm{C} _ {i}) \rightarrow \mathrm{H} \left(\bigoplus_ {i \in \mathbf {I}} \mathrm{C} _ {i}\right), \quad \mathrm{H} \left(\prod_ {i \in \mathbf {I}} \mathrm{C} _ {i}\right)\rightarrow \prod_ {i \in \mathbf {I}} \mathrm{H} (\mathrm{C} _ {i})\tag{2}
\]

are bijective.

For every filtered inductive system of complexes \((\mathbf{C}_{i})_{i\in\mathbf{I}}\) , the canonical homomorphism

\[
\varinjlim_ {i \in \mathrm{I}} \mathrm{H} (\mathrm{C} _ {i}) \to \mathrm{H} (\varinjlim_ {i \in \mathrm{I}} \mathrm{C} _ {i})\tag{3}
\]

is bijective.

This follows immediately from II, p.~14, cor. 1 to prop. 7, p.~11, cor. to prop. 5, and p.~91, prop. 3.

